\section{Open-World Planning and Learning}
\label{sec:rw_planning_learning}

Planning and learning provide the reasoning and decision-making capabilities that allow agents to select actions in pursuit of goals. This section reviews the major paradigms and identifies the limitations that the contributions of this dissertation address.

\paragraph{Classical symbolic planning.}
Classical symbolic planning, formalized through STRIPS \citep{fikes1971strips} and PDDL \citep{mcdermott1998pddl}, enables structured long-horizon reasoning by searching over sequences of operators with well-defined preconditions and effects. Modern planners such as Fast Downward \citep{helmert2006fast} achieve strong performance across diverse domains through domain-independent heuristics. However, symbolic planning assumes that the domain model is complete and correct; when the model is inaccurate or missing operators, generated plans may fail during execution. This brittleness under model mismatch is precisely the problem addressed by the hybrid planning--learning framework in \cref{ch:rapid_learn}, where symbolic plans provide structure for targeted adaptation rather than being treated as infallible specifications.

\paragraph{Reinforcement learning.}
Reinforcement learning offers a complementary approach, enabling agents to learn effective behaviors through interaction without requiring explicit models. Deep RL algorithms, including DQN \citep{mnih2015human}, PPO \citep{schulman2017proximal}, and SAC \citep{haarnoja2018soft}, have demonstrated success in complex control tasks. Model-based approaches, including world models such as Dreamer \citep{hafner2023mastering} and planning-based methods such as MuZero \citep{schrittwieser2020mastering}, learn environment dynamics and use them for decision-making. These methods reduce the need for hand-specified models but typically assume that the learned dynamics remain stationary---an assumption that open-world settings violate. The adaptation methods in \cref{ch:rapid_learn} and the generalization methods in \cref{ch:generalization} both build on reinforcement learning but address different aspects of this stationarity limitation: the former enables targeted relearning when dynamics change, while the latter produces representations that are invariant to certain classes of variation from the outset.

\paragraph{Task and motion planning.}
Task and motion planning (TAMP) systems \citep{garrett2021integrated, kaelbling2011hierarchical} integrate symbolic task planning with continuous motion planning, enabling agents to reason jointly about high-level task structure and low-level geometric feasibility. These systems represent a sophisticated form of the hybrid planning--learning framework, but they typically assume complete and correct domain models at both levels. When the environment changes in ways that invalidate either the task-level operators or the motion-level feasibility assumptions, TAMP systems lack mechanisms for acquiring new behaviors. The approach in \cref{ch:rapid_learn} can be viewed as extending the hybrid paradigm with the capacity for targeted learning when model assumptions fail.

\paragraph{Transfer and continual learning.}
When an agent encounters a changed environment, transfer learning and continual learning provide mechanisms for leveraging prior experience. Policy reuse \citep{glatt2017policy}, progressive neural networks \citep{rusu2016progressive}, and elastic weight consolidation \citep{kirkpatrick2017overcoming} address different aspects of this problem. However, these methods treat adaptation as a monolithic learning problem over the full state-action space: they do not exploit the compositional structure of tasks to focus learning on the specific component that failed. As the experimental results in \cref{ch:rapid_learn} demonstrate, this global approach is substantially less sample-efficient than the targeted adaptation enabled by action abstractions when novelty affects only a specific component of the task.

\paragraph{Open-world novelty handling.}
The DARPA SAIL-ON program \citep{senator2019science} catalyzed research specifically targeting agents that must detect, characterize, and accommodate novelty. This program motivated formal definitions of open-world novelty \citep{boult2021towards, langley2020open} and agent architectures for novelty handling \citep{muhammad2021novelty, goel2024neurosymbolic}. Several approaches explored how agents can adapt their policies or models in response to detected novelty \citep{balloch2022novgrid, balloch2023neuro}, and neuro-symbolic world models have been proposed for updating internal representations under novelty. The benchmarks in \cref{ch:benchmarks} and the adaptation framework in \cref{ch:rapid_learn} build directly on this line of work, providing both the evaluation infrastructure and the algorithmic machinery for systematic novelty handling.

\paragraph{Connection to this dissertation.}
Across these paradigms, a recurring limitation is the difficulty of reconciling structured reasoning (which provides compositional decomposition but assumes correct models) with learned behavior (which can acquire new capabilities but lacks structural guidance for efficient adaptation). The contributions of this dissertation address this gap through structured abstractions: action abstractions that bridge symbolic failure diagnosis with targeted learning (\cref{ch:rapid_learn}), and interaction abstractions that organize learned behavior around task-relevant invariants rather than embodiment-specific details (\cref{ch:generalization}). The benchmarks in \cref{ch:benchmarks} provide the controlled evaluation infrastructure needed to measure progress on these problems.